\documentclass[UTF8,zihao=-4]{ctexart}
\usepackage[a4paper,margin=1.9cm]{geometry}
\usepackage{hyperref}
\usepackage{array}
\usepackage{longtable}
\usepackage{verbatim}
\usepackage{xcolor}
\hypersetup{hidelinks}
\setlength{\parindent}{0pt}
\setlength{\parskip}{0.45em}
\emergencystretch=4em
\sloppy
\title{00-2025风格零基础期末总复习资料}
\author{}
\date{}
\begin{document}
\maketitle
\setcounter{tocdepth}{1}
\tableofcontents
\newpage
\section{编译原理 2025 风格零基础期末突击总复习资料}


\subsection{使用方法}


这份资料按 Slides/Slides/handout 与三份 Recap PPT 重排复习重点。新老师 PPT 覆盖前端、IR/SSA、中端数据流/符号执行/指针分析和后端代码生成。复习时按下面顺序：


\begin{quote}
\footnotesize
\begin{verbatim}
词法分析 -> LL 语法分析 -> IR/SSA -> 数据流分析 -> 符号执行 -> 过程间/指针分析 -> 目标代码/寄存器分配/指令调度
\end{verbatim}
\end{quote}


旧 Review/compile.pdf 中的 LR、ANTLR、SDD、TAC 不再作为模拟卷主线；其中 TAC/SDD 只在 IR 生成相关 PPT 范围内保留。


\subsection{一、Slides 主线题型蓝图}


\begin{quote}
\footnotesize
\begin{verbatim}
1. 词法分析：正则表达式 -> Thompson NFA -> 子集构造 DFA -> DFA 最小化。
2. LL 语法分析：消除左递归 -> FIRST -> FOLLOW -> 预测分析表 -> LL(1) 判断。
3. IR/SSA：三地址码、基本块、CFG、支配关系、支配边界、phi 插入与 phi 转出。
4. 数据流分析：方向、meet、transfer、gen/kill、worklist、IN/OUT。
5. 符号执行：敏感性定义、CNF、Tseitin、同时可满足证明。
6. 过程间/指针分析：调用上下文、context-sensitive、Andersen 约束、points-to 集合。
7. 后端：目标代码、地址模式、冲突图、寄存器分配、spill、指令调度。
\end{verbatim}
\end{quote}


\subsection{二、词法分析模板}


\subsubsection{Thompson NFA}


\begin{quote}
\footnotesize
\begin{verbatim}
字符 a：start --a--> end
空串 epsilon：start --epsilon--> end
选择 A|B：新 start epsilon 到 A/B 起点，A/B 终点 epsilon 到新 end
连接 AB：A 终点 epsilon 到 B 起点
闭包 A*：新 start epsilon 到 A 起点和新 end，A 终点 epsilon 到 A 起点和新 end
\end{verbatim}
\end{quote}


答题步骤：


\begin{quote}
\footnotesize
\begin{verbatim}
1. 按正则表达式括号从内到外拆分。
2. 每个子表达式套 Thompson 模板。
3. 给 NFA 状态编号。
4. 起点画入箭头，终点画双圈。
\end{verbatim}
\end{quote}


\subsubsection{子集构造}


\begin{quote}
\footnotesize
\begin{verbatim}
epsilon-closure(S)：从 S 出发只走 epsilon 边到达的状态集合。
move(S,a)：从 S 出发读 a 到达的状态集合。
Start = epsilon-closure({NFA start})。
DFA 终态：NFA 状态集合中含 NFA 终态。
\end{verbatim}
\end{quote}


表格写法：


\begin{quote}
\footnotesize
\begin{verbatim}
DFA状态 | NFA状态集合 | on a | on b | 终态
A       | {...}       | B    | C    | 否
B       | {...}       | B    | D    | 是
\end{verbatim}
\end{quote}


\subsubsection{DFA 最小化}


\begin{quote}
\footnotesize
\begin{verbatim}
P0 = {终态集合, 非终态集合}
按每个输入字符的转移目标所在组拆分。
重复拆分直到分组不再变化。
每个分组变成一个最小 DFA 状态。
\end{verbatim}
\end{quote}


\subsection{三、LL 语法分析模板}


直接左递归：


\begin{quote}
\footnotesize
\begin{verbatim}
A -> A alpha | beta
\end{verbatim}
\end{quote}


消除：


\begin{quote}
\footnotesize
\begin{verbatim}
A  -> beta A'
A' -> alpha A' | epsilon
\end{verbatim}
\end{quote}


FIRST：


\begin{quote}
\footnotesize
\begin{verbatim}
FIRST(a) = {a}
FIRST(epsilon) = {epsilon}
FIRST(XY)：先看 FIRST(X)，若 X 能推出 epsilon，再看 FIRST(Y)。
\end{verbatim}
\end{quote}


FOLLOW：


\begin{quote}
\footnotesize
\begin{verbatim}
开始符号 S：$ in FOLLOW(S)
A -> alpha B beta：FIRST(beta)-{epsilon} 加入 FOLLOW(B)
A -> alpha B 或 beta 能推出 epsilon：FOLLOW(A) 加入 FOLLOW(B)
\end{verbatim}
\end{quote}


预测表：


\begin{quote}
\footnotesize
\begin{verbatim}
对 A -> alpha：
FIRST(alpha)-{epsilon} 中每个 a 填 M[A,a]
若 epsilon in FIRST(alpha)，FOLLOW(A) 中每个 b 填 M[A,b]
表中无多重入口，则是 LL(1)。
\end{verbatim}
\end{quote}


2025 原题模板：


\begin{quote}
\footnotesize
\begin{verbatim}
A -> B C
B -> B P | r
C -> t | u
\end{verbatim}
\end{quote}


消除左递归：


\begin{quote}
\footnotesize
\begin{verbatim}
A  -> B C
B  -> r B'
B' -> P B' | epsilon
C  -> t | u
\end{verbatim}
\end{quote}


FIRST 与 FOLLOW：


\begin{quote}
\footnotesize
\begin{verbatim}
FIRST(A)={r}
FIRST(B)={r}
FIRST(B')={P,epsilon}
FIRST(C)={t,u}
FOLLOW(A)={$}
FOLLOW(B)={t,u}
FOLLOW(B')={t,u}
FOLLOW(C)={$}
\end{verbatim}
\end{quote}


预测表：


\begin{quote}
\footnotesize
\begin{verbatim}
M[A,r]  = A -> B C
M[B,r]  = B -> r B'
M[B',P] = B' -> P B'
M[B',t] = B' -> epsilon
M[B',u] = B' -> epsilon
M[C,t]  = C -> t
M[C,u]  = C -> u
\end{verbatim}
\end{quote}


\subsection{四、指令调度模板}


目的：


\begin{quote}
\footnotesize
\begin{verbatim}
在不改变程序语义的前提下重排指令，减少流水线停顿，提高并行执行效率。
\end{verbatim}
\end{quote}


三类依赖：


\begin{quote}
\footnotesize
\begin{verbatim}
True dependence / RAW：i 写 x，j 读 x，i 必须在 j 前。
Anti dependence / WAR：i 读 x，j 写 x，j 不能提前破坏 i 读到的旧值。
Output dependence / WAW：i 写 x，j 写 x，二者顺序影响最终写入值。
\end{verbatim}
\end{quote}


寄存器重命名：


\begin{quote}
\footnotesize
\begin{verbatim}
能消除 WAR 和 WAW。
不能消除 RAW。
\end{verbatim}
\end{quote}


答题步骤：


\begin{quote}
\footnotesize
\begin{verbatim}
1. 每条指令一个节点。
2. RAW 画实线边。
3. WAR 和 WAW 画虚线边。
4. 对写后写、读后写冲突做寄存器重命名。
5. 对剩余依赖图给一个拓扑序。
\end{verbatim}
\end{quote}


\subsection{五、支配问题模板}


定义：


\begin{quote}
\footnotesize
\begin{verbatim}
Dom(n)：从入口到 n 的每条路径都经过的节点集合。
PostDom(n)：从 n 到出口的每条路径都经过的节点集合。
idom(n)：n 的最近严格支配者。
Dom Frontier(n)：n 支配某个前驱，但不严格支配该节点的位置。
\end{verbatim}
\end{quote}


计算 Dom：


\begin{quote}
\footnotesize
\begin{verbatim}
Dom(entry) = {entry}
其他节点 Dom 初值为全集
Dom(n) = {n} union intersection(Dom(p) for p in pred(n))
迭代到不再变化
\end{verbatim}
\end{quote}


SSA 中 phi 放置：


\begin{quote}
\footnotesize
\begin{verbatim}
变量在多个基本块中被定义时，需要看这些定义块的支配边界。
phi 放在多条控制流汇合且不同定义可能到达的位置。
更精确地说：在定义块的迭代支配边界处插入该变量的 phi。
\end{verbatim}
\end{quote}


\subsection{六、符号执行与 SAT 模板}


敏感性：


\begin{quote}
\footnotesize
\begin{verbatim}
Path-sensitive：区分不同执行路径。
Flow-sensitive：区分程序点和语句顺序。
Context-sensitive：区分函数调用上下文。
Field-sensitive：区分对象字段。
\end{verbatim}
\end{quote}


CNF：


\begin{quote}
\footnotesize
\begin{verbatim}
CNF 是子句的合取，每个子句是文字的析取。
(a or not b) and (c or d)
\end{verbatim}
\end{quote}


Tseitin：


\begin{quote}
\footnotesize
\begin{verbatim}
1. 给每个子表达式引入新变量。
2. 为新变量和子表达式的语义关系写 CNF。
3. 加入根变量为真的子句。
4. 得到与原式同时可满足的 CNF。
\end{verbatim}
\end{quote}


常用子句：


\begin{quote}
\footnotesize
\begin{verbatim}
x <-> (a and b)
(not x or a) and (not x or b) and (x or not a or not b)

x <-> (a or b)
(x or not a) and (x or not b) and (not x or a or b)

x <-> not a
(not x or not a) and (x or a)
\end{verbatim}
\end{quote}


同时可满足证明模板：


\begin{quote}
\footnotesize
\begin{verbatim}
F1 = (a0 or ... or an or c) and (b0 or ... or bm or not c)
F2 = (a0 or ... or an or b0 or ... or bm)

F1 -> F2：
若 c=false，则第一子句要求某个 ai=true。
若 c=true，则第二子句要求某个 bj=true。
因此 F2=true。

F2 -> F1：
若某个 ai=true，令 c=false。
若某个 bj=true，令 c=true。
因此 F1=true。
\end{verbatim}
\end{quote}


\subsection{七、Andersen 指针分析模板}


points-to 集合：


\begin{quote}
\footnotesize
\begin{verbatim}
pts(p) = p 能指向的对象集合
\end{verbatim}
\end{quote}


四类语句：


\begin{quote}
\footnotesize
\begin{verbatim}
y = &x：x in pts(y)
y = x：pts(x) subset pts(y)
*y = x：for v in pts(y), pts(x) subset pts(v)
y = *x：for v in pts(x), pts(v) subset pts(y)
\end{verbatim}
\end{quote}


闭包步骤：


\begin{quote}
\footnotesize
\begin{verbatim}
1. 先处理取地址语句，得到初始 points-to 集。
2. 处理复制语句，传播集合包含关系。
3. 处理 store 语句，把右侧集合写入左侧指向对象的集合。
4. 处理 load 语句，把被指对象的集合读入左侧变量。
5. 有新增元素就继续传播，直到不再变化。
\end{verbatim}
\end{quote}


\subsection{八、PPT 新增重点}


IR/SSA：


\begin{quote}
\footnotesize
\begin{verbatim}
三地址码：每条语句最多一个主要运算。
基本块：只有入口第一条和出口最后一条会改变控制流。
CFG：节点是基本块，边是可能的控制流。
SSA：每个变量只有一个定义，汇合处用 phi 表示来源选择。
\end{verbatim}
\end{quote}


数据流分析：


\begin{quote}
\footnotesize
\begin{verbatim}
先判断方向：前向或后向。
再判断性质：may 用 union，must 用 intersection。
transfer 常写成 OUT = gen union (IN - kill)。
worklist 到所有 IN/OUT 不再变化时停止。
\end{verbatim}
\end{quote}


过程间分析：


\begin{quote}
\footnotesize
\begin{verbatim}
context-insensitive 会合并不同调用点的信息。
context-sensitive 会区分调用上下文。
1-call-string 用最近一个调用点作为上下文。
精度提高通常会带来更高分析开销。
\end{verbatim}
\end{quote}


目标代码与寄存器分配：


\begin{quote}
\footnotesize
\begin{verbatim}
目标代码生成要说明 load、store、calculation 和地址模式。
冲突图中，同一程序点同时活跃的变量之间有边。
k 个寄存器对应 k 着色问题。
如果图无法 k 着色，就需要 spill。
\end{verbatim}
\end{quote}


\subsection{九、最后背诵清单}


\begin{quote}
\footnotesize
\begin{verbatim}
Thompson 五模板
子集构造三操作：epsilon-closure、move、终态判断
DFA 最小化初分组和拆分
左递归消除公式
FIRST/FOLLOW/预测表规则
RAW/WAR/WAW 定义
寄存器重命名能消除 WAR/WAW
Dom/PostDom/Dom Frontier 定义
支配边界用于 phi 放置
数据流方向、meet、transfer、gen/kill
Path/Flow/Context/Field sensitivity
CNF 与 Tseitin 三类子句
Andersen 四类约束
目标代码、冲突图、spill、指令调度
\end{verbatim}
\end{quote}

\end{document}
